\section{Conclusion}
\label{sec:conclusion}

We introduced calibration analysis to copy mechanisms, asking whether a specialized copy-decision model produces better-calibrated probabilities than a joint pointer-generator. Using a synthetic task with ground-truth copy labels, we find that both models are well-calibrated (ECE $< 0.01$), with no significant difference between them. However, specialization produces substantially more accurate copy decisions (F1: 0.985 vs.\ 0.969) and better Brier scores (0.013 vs.\ 0.024). The practical takeaway is clear: if you need reliable copy decisions, train a specialized model---the accuracy gains are large even though calibration is comparable.

Our most surprising finding is that joint generation training appears to implicitly regularize calibration, an effect that strengthens with model capacity. This suggests a broader principle: auxiliary tasks that are harder than the target decision may improve calibration of that decision even when they do not share a direct training signal.

Future work should validate these findings on real-world copy tasks such as CNN/Daily Mail summarization, explore end-to-end specialized models with their own encoders, and test whether the implicit calibration regularization effect holds at the scale of modern pre-trained models like BART and T5.